(1) 单色物光为缓慢变化的复振幅 $O(x)$，参考光为 $R_0e^{ik_rx}$，其中 $R_0>0$。探测器只测强度。写出全息干涉图 $I(x)$ 并分解为零级和两个载频交叉项。

(2) 设物光空间频谱仅在 $|k|<B$ 内，给出足够大的 $|k_r|$ 以使两个交叉侧带与直流项分开，解释如何在频域重建物光复振幅。

(3) 对局部近似常幅物光 $O=Ae^{i\phi}$ 求强度条纹可见度及最优参考强度，说明两束不匹配时可见度的变化。

(4) 讨论参考光复相位未知、探测有限像素间距与散斑相干性对重建的限制，并解释所谓双像或共轭项的来源。
